% !TeX program = xelatex
\documentclass[UTF8,fontset=fandol,zihao=5,a4paper]{ctexart}
\usepackage[margin=20mm]{geometry}
\usepackage{amsmath,amssymb,booktabs,tabularx,array,graphicx}
\usepackage{tikz}
\usetikzlibrary{arrows.meta,positioning,fit,backgrounds}
\usepackage[hidelinks]{hyperref}
\definecolor{styleblue}{HTML}{397BAC}
\definecolor{contentgreen}{HTML}{39816E}
\definecolor{floworange}{HTML}{B97C31}
\definecolor{timepurple}{HTML}{8864AB}
\tikzset{
  box/.style={draw=black!45,fill=white,rounded corners=3pt,align=center,
    font=\small,inner sep=6pt,minimum width=3cm,minimum height=.8cm},
  stylebox/.style={box,draw=styleblue,fill=styleblue!5},
  contentbox/.style={box,draw=contentgreen,fill=contentgreen!5},
  flowbox/.style={box,draw=floworange,fill=floworange!5},
  timebox/.style={box,draw=timepurple,fill=timepurple!5},
  arrow/.style={-{Latex[length=2mm]},line width=.8pt},
  condition/.style={arrow,dashed},
}
\newcommand{\code}[1]{\texttt{\detokenize{#1}}}
\newcommand{\LN}{\operatorname{LN}}
\newcommand{\MHA}{\operatorname{MHA}}
\newcommand{\FFN}{\operatorname{FFN}}
\setlength{\parindent}{0pt}
\setlength{\parskip}{5pt}
\linespread{1.12}
\ctexset{section={format=\large\bfseries,beforeskip=12pt,afterskip=6pt},
  subsection={format=\normalsize\bfseries,beforeskip=9pt,afterskip=4pt}}

\begin{document}
\begin{center}
{\LARGE\bfseries Glyph2Hand}\par
\vspace{4pt}
{\large 纯噪声 flow：独立内容编码与多参考风格条件}\par
\vspace{3pt}
{\small 2026 年 10 月 10 日\quad 当前架构与训练实现，检查点版本 3}
\end{center}

\section{任务与当前实现}

输入为同一作者的 $N$ 个参考手写字 $\{H_i\}_{i=1}^N$ 和目标标准字 $K_\star$。
参考字不再配同字标准字。模型从标准高斯噪声生成目标字的手写图，目标标准字通过独立内容 token 提供条件。
目标手写真值 $H_\star$ 仅用于训练状态和速度监督，不放入参考集合。
实际标准字体由数据准备参数 \code{--font} 指定，训练说明使用文津宋体；代码不硬编码字体。

风格分支将每个参考字独立浓缩为 $Z$ 个描述，然后直接拼接。
内容分支采用相同的 patch/SA 结构，参数与 SA 深度独立，不做最后的 CA 浓缩。
主干堆叠 \textbf{SA $\to$ Content CA $\to$ Style CA $\to$ FFN}。
训练使用 flow MSE 和作者风格损失，每个 batch 更新一次；验证、预览使用当前模型。
完整配置、运行命令、日志与 I/O 处理见上一级 \code{README.md} 和 \code{TRAINING.md}。
\code{README.v1-with-diagrams.*} 为历史设计；本文描述代码行为，不提供生成质量结论。

\subsection{张量与符号}

采用通道在前布局。批量为 $B$，通道 $C=1$，图像边长 $n$，patch 边长 $p$，
空间 token 数 $L=(n/p)^2$，特征维 $d$，每字描述数 $Z=\code{description_tokens}$。
\begin{center}\small
\begin{tabularx}{\linewidth}{@{}l l X@{}}
\toprule
张量 & 形状 & 用途 \\
\midrule
参考手写字 & $[B,N,C,n,n]$ & 同一 episode 来自同一作者 \\
目标标准字 & $[B,C,n,n]$ & 内容条件，独立编码 \\
参考有效掩码 & $[B,N]$ & True 表示有效，每个样本至少一个 \\
内容 $C_\star$ & $[B,L,d]$ & 保留空间位置的内容序列 \\
风格 $U$ & $[B,NZ,d]$ & 各参考描述直接拼接 \\
当前状态 / 速度 & $[B,C,n,n]$ & 像素空间 flow \\
时间 & $[B]$ & 连续流位置，训练均匀采样 \\
\bottomrule
\end{tabularx}\end{center}
字图映射为 $[-1,1]$，白色 $+1$、黑色 $-1$；噪声、中间状态和速度不受该范围约束。
底层模型可配置通道数，现有数据处理、训练器和预览仅支持单通道。

\clearpage
\section*{架构图：三条 patch 分支与两类条件读取}
\begin{center}
\resizebox{!}{208mm}{%
\begin{tikzpicture}[x=.95cm,y=1cm]
\path[use as bounding box] (0,1) rectangle (16,-18.5);
\node[stylebox] (ref) at (2,0) {同作者参考手写字\\$H_1,\ldots,H_N$};
\node[contentbox] (standard) at (14,0) {目标标准字\\$K_\star$};
\node[flowbox] (noise) at (8,0) {纯高斯起点\\$x_0=\epsilon\sim\mathcal N(0,I)$};
\node[stylebox] (refpatch) at (2,-1.6) {patch 投影 $Cp^2\to d$\\加一次二维位置 $P_{xy}$};
\node[contentbox] (contentpatch) at (14,-1.6) {patch 投影 $Cp^2\to d$\\加一次二维位置 $P_{xy}$};
\node[flowbox] (state) at (8,-1.6) {当前状态 $x_t$};
\node[stylebox] (refsa) at (2,-3.2) {SA $\to$ FFN\\$\times B_s$；逐字共享参数};
\node[contentbox] (contentsa) at (14,-3.2) {SA $\to$ FFN\\$\times B_c$；独立参数与深度};
\node[flowbox] (genpatch) at (8,-3.2) {patch 投影 $Cp^2\to d$\\加一次二维位置 $P_{xy}$};
\node[stylebox] (pool) at (2,-4.9) {可学习查询 $Q_Z$\\一层 CA $\to$ FFN\\查询与输出不加位置编码};
\node[contentbox] (content) at (14,-4.9) {内容 $C_\star$：$L\times d$\\保留全部空间 token\\不做 CA 浓缩};
\node[stylebox] (style) at (2,-6.7) {每字 $Z$ 个描述直接拼接\\$U$：$NZ\times d$\\无参考集合浓缩};
\node[flowbox] (sa) at (8,-7.4) {SA\\AdaLN-Zero + 残差门};
\node[flowbox] (cca) at (8,-9.2) {Content CA\\普通 LN + 残差};
\node[flowbox] (sca) at (8,-11) {Style CA\\普通 LN + 残差\\屏蔽补齐描述};
\node[flowbox] (ffn) at (8,-12.8) {FFN\\AdaLN-Zero + 残差门};
\node[timebox] (time) at (14,-12.8) {$t\to$ sin/cos $\to$ MLP\\各块独立 $6d$ 时间调制\\作用于 SA / FFN};
\node[flowbox] (head) at (8,-14.8) {时间调节 LN + 共享输出投影\\$d\to Cp^2$，按位置拼回速度};
\node[flowbox] (euler) at (8,-16.5) {$x_{t+\Delta t}=x_t+\Delta t\,v_\theta$\\到终点后裁剪 $[-1,1]$};
\foreach \a/\b in {ref/refpatch,refpatch/refsa,refsa/pool,pool/style,
standard/contentpatch,contentpatch/contentsa,contentsa/content,
noise/state,state/genpatch,genpatch/sa,sa/cca,cca/sca,sca/ffn,ffn/head,head/euler}{
  \draw[arrow] (\a.south)--(\b.north);
}
\draw[condition,draw=contentgreen] (content.south)--(14,-9.2)--(cca.east);
\draw[condition,draw=styleblue] (style.south)--(2,-11)--(sca.west);
\draw[condition,draw=timepurple] (time.north)--(11.2,-12)--(11.2,-7.4)--(sa.east);
\draw[condition,draw=timepurple] (time.west)--(ffn.east);
\draw[condition,draw=timepurple] (time.south)--(14,-14.8)--(head.east);
\draw[arrow,draw=floworange] (euler.west)--(5,-16.5)--(5,-1.6)--(state.west);
\node[font=\scriptsize,anchor=east,align=right,text=floworange] at (4.8,-8.5) {未到终点\\继续积分};
\begin{scope}[on background layer]
\node[draw=floworange!50,fill=floworange!2,rounded corners=4pt,
fit=(sa)(cca)(sca)(ffn),inner sep=10pt] {};
\end{scope}
\node[font=\small,text=floworange] at (11.7,-6.4) {主干块 $\times B_g$};
\node[font=\small,align=center] at (8,-18) {条件编码在积分前各计算一次；每层分别读取同一份 $C_\star$ 与 $U$};
\end{tikzpicture}}
\end{center}
{\small 风格、内容和状态均在 patch 投影后加一次位置编码；风格浓缩查询及描述不加位置编码。
图中的 CA 表示主干独立注意力子层；风格浓缩模块含一层 CA 和一个 FFN，均有前置 LN 与残差。}

\clearpage
\section{风格与内容编码}

对第 $i$ 个参考手写字，定义
\[
R_i=E_s\bigl(\operatorname{Patch}_s(H_i)+P_{xy}\bigr)\in\mathbb R^{L\times d},
\qquad U_i=\operatorname{Pool}(Q_Z,R_i)\in\mathbb R^{Z\times d}.
\]
每字独立编码，各字共享 patch 投影、$B_s=\code{local_depth}$ 个 SA--FFN 块和可学习查询 $Q_Z$。
每个 SA--FFN 子层采用普通前置 LN 与残差，FFN 为两层线性映射，中间 GELU。
\textbf{Pool 中恰好只有一层 CA}，随后一个 FFN：
\begin{align*}
\widetilde Q&=Q_Z+\MHA(\LN_q(Q_Z),\LN_r(R_i),\LN_r(R_i)),\\
U_i&=\widetilde Q+\FFN(\LN_f(\widetilde Q)).
\end{align*}
$Q_Z$ 是共享训练参数，以标准差 $0.02$ 的正态分布初始化；查询和浓缩输出均不添加位置编码。
浓缩后不再做自注意力，不保留原始空间 token 旁路。
直接拼接 $U=[U_1;\ldots;U_N]$，主干读取全部 $NZ$ 个描述，没有集合浓缩层或其他汇总查询。
描述槽位不等于预先指定的笔画类别，内容与风格的严格解耦也不由结构保证。

内容分支为
\[
C_\star=E_c\bigl(\operatorname{Patch}_c(K_\star)+P_{xy}\bigr)\in\mathbb R^{L\times d}.
\]
其 patch 投影和 SA--FFN 与单参考字风格编码同构，参数完全独立，
$B_c=\code{content_depth}$ 可单独调整；不执行最后的 Pool，保持空间 token。
内容和风格都不接受时间/当前状态输入，可在积分步之间缓存，训练更新后缓存需重算。

\subsection{参考重排与掩码}

参考重排只会重排 $U$ 的整组行。主干 Style CA 对 K/V 同步重排不变，因此同样重排有效掩码后，
固定内容、状态和时间的速度结果不变，允许浮点误差。
代码不对拼接序列再添加组序号或位置编码。
参考掩码 True 表示有效，扩展到每字 $Z$ 个描述后，在每层 Style CA 中取反为 padding mask。
省略掩码时全部有效，拒绝一个样本所有参考均无效的情况。
补齐图像必须使用有限数值；数据集补白 $+1$，补齐槽仍会参加逐字编码。

\section{主干与时间调制}

主干仅将当前流状态分块：$T=\operatorname{Patch}_g(x_t)+P_{xy}$。
目标标准字只作为内容 memory，不拼为图像通道，也不参与起点构造。
每个 block 使用独立时间映射生成 SA/FFN 的缩放、平移和残差门：
\[
(a_s,b_s,g_s,a_f,b_f,g_f)=\operatorname{Split}_6(\operatorname{Linear}(\operatorname{SiLU}(e(t)))),
\quad D_j(T)=(1+a_j)\odot\LN_j(T)+b_j.
\]
六个 $d$ 维向量在该样本全部 token 间广播，gate 不使用 sigmoid。
完整顺序为
\begin{align*}
A&=D_s(T),\\
T_s&=T+g_s\odot\MHA_s(A,A,A),\\
T_c&=T_s+\MHA_c(\LN_{cq}(T_s),\LN_{ck}(C_\star),\LN_{ck}(C_\star)),\\
T_u&=T_c+\MHA_u(\LN_{uq}(T_c),\LN_{uk}(U),\LN_{uk}(U)),\\
T^+&=T_u+g_f\odot\FFN(D_f(T_u)).
\end{align*}
Content CA 和 Style CA 参数独立，均为普通前置 LN 和残差，没有单独的时间调节或零门。
SA 和 FFN 的 LN 无仿射，$\varepsilon=10^{-6}$；普通 LN 有可学习仿射，$\varepsilon=10^{-5}$。
所有 MHA 为 batch-first、dropout $0$、带 bias，Q/K/V 独立投影，不绑定 K/V，调用不返回注意力权重。

时间特征频率为 $f_k=\exp(-\log(10000)k/(F/2))$，
原始特征拼接 $\cos(t\tau f_k)$、$\sin(t\tau f_k)$，再经 Linear--SiLU--Linear 得到 $e(t)$。
$F$ 为时间特征维，$\tau$ 为时间放大倍数。二维位置按行/列各占 $d/2$ 维，每轴拼接 sin/cos。
位置和频率为非持久化 buffer，不写入 state dict。

输出头用独立 SiLU--Linear 产生 $2d$ 个 scale/shift，调节无仿射 LN，再共享线性投影为每块 $Cp^2$ 个速度并拼回图像。
主干 $6d$ 调制、输出头 $2d$ 调制和速度投影均零初始化；初始速度为零，普通 CA 仍有残差，整个 block 不保证恒等。

\section{双损失与作者配对方式}

\subsection{纯噪声 flow MSE}

独立采样 $\epsilon\sim\mathcal N(0,I)$ 和 $t\sim\mathcal U(0,1)$：
\begin{align*}
x_0&=\epsilon,\qquad x_t=(1-t)\epsilon+tH_\star,\qquad u=H_\star-\epsilon,\\
\mathcal L_{\rm flow}&=\mathbb E\left[\frac{1}{Cn^2}\|v_\theta(x_t,t,C_\star,U)-u\|_F^2\right].
\end{align*}
训练器用 float32 平均样本、通道和像素的平方误差，不裁剪训练状态。
标准字只进入内容编码，参考字只进入风格编码；两者与主干共同反向训练。

\subsection{作者风格损失}

将 $U$ 整理为 $[B,N,Z,d]$。作者损失先对有效参考的同描述槽取均值：
\[
A_{i,z}=\frac{\sum_r m_{i,r}U_{i,r,z}}{\sum_r m_{i,r}},\qquad
\bar A_{i,z}=\frac{A_{i,z}}{\max(\|A_{i,z}\|_2,10^{-8})},\qquad
c_{ijz}=\bar A_{i,z}^{\top}\bar A_{j,z}.
\]
\textbf{这个均值只用于计算损失，主干依然读取全部 $NZ$ 个描述。}
没有为作者损失增加可学习集合浓缩层。均值、归一化及相似度比较均在关闭 autocast 的 float32 上下文计算。
设作者索引为 $a_i$，取当前 batch 的有序配对
$P_s=\{(i,j):i\ne j,a_i=a_j\}$、$P_d=\{(i,j):a_i\ne a_j\}$：
\begin{align*}
\mathcal L_s&=\frac{1}{|P_s|Z}\sum_{(i,j)\in P_s}\sum_z(1-c_{ijz}),\\
\mathcal L_d&=\frac{1}{|P_d|Z}\sum_{(i,j)\in P_d}\sum_z c_{ijz}^2,\\
\mathcal L_{\rm author}&=\mathcal L_s+\mathcal L_d,\qquad
\mathcal L_{\rm train}=\mathcal L_{\rm flow}+0.05\,\mathcal L_{\rm author}\quad\text{（默认）}.
\end{align*}
两类分别平均，再相加；缺少某类时对应项为保持梯度连接的零。单样本 batch 作者项为零。
函数支持两类独立非负权重，默认均为 $1$；总作者项权重可在训练配置修改，$0$ 可关闭。
同作者鼓励一致，异作者鼓励正交，相反方向不是零损失解。
训练与验证均统计当前模型的双损失总目标，日志为 \code{train_loss} 和 \code{val_loss}。

\subsection{为什么没有二次方、四次方的数据枚举}

修改前的数据采样：每个 episode 独立均匀选作者，再从作者无放回抽 $N+1$ 个不同字符，
首个为目标，余下为参考。旧作者损失比较集合浓缩后的固定槽，且只在当次 batch 内配对，
没有跨 batch 缓存、专门保证同作者重复的采样，或对全部参考集合做组合枚举。
修改后保留相同采样方式，只改为上述同槽均值表示。

一个 batch 的作者出现次数为 $k_a$ 时，同作者有序对数为
$\sum_a k_a(k_a-1)$，异作者为 $B(B-1)-\sum_a k_a(k_a-1)$。
例如作者 $[a,a,b,c]$ 有 $2$ 个同作者、$10$ 个异作者有序对；两类各自取平均。
均匀选择 $W$ 个作者时，同作者有序对的期望为 $B(B-1)/W$，一批完全没有同作者配对是可能的。
全量枚举不同参考集合的配对空间可以很大，但代码只比较已抽到的 $B$ 个 episode，
两类相似度计算都是 $O(B^2Zd)$，没有四次方枚举。异作者公式中的平方是余弦值平方，不是采样次数平方。

\section{配置与采样}

\begin{center}\small
\begin{tabularx}{\linewidth}{@{}l l X@{}}
\toprule
模型字段 & 默认值 & 含义 \\
\midrule
\code{image_size / patch_size / in_channels} & $64/8/1$ & $L=64$，单通道 \\
\code{hidden_size / num_heads / mlp_ratio} & $256/8/2$ & 每头 $32$，FFN 内宽 $512$ \\
\code{local_depth / content_depth} & $8/8$ & 风格、内容 SA 深度独立 \\
\code{generator_depth} & $8$ & 双 CA 主干深度 \\
\code{description_tokens} & $24$ & 每字 $Z$ 个描述 \\
\code{time_frequency_size / time_scale} & $256/1000$ & 时间编码维度与倍数 \\
\bottomrule
\end{tabularx}\end{center}
图像边长须整除 patch 边长，$d$ 可被头数和 $4$ 整除，时间特征维为正偶数，三个 SA 深度允许为 $0$。
时间特征维默认在类定义时确定为 $256$，单独改 $d$ 不会自动同步。
默认参数总数 $20{,}890{,}688$：参考 $4{,}767{,}232$，内容 $4{,}233{,}472$，主干 $11{,}889{,}984$。

\begin{center}\small
\begin{tabularx}{\linewidth}{@{}l l X@{}}
\toprule
训练字段 & 默认值 & 含义 \\
\midrule
\code{steps / batch_size} & $100000/64$ & 每 batch 一次优化器更新 \\
\code{min_references / max_references} & $4/4$ & 默认 $NZ=96$；可设数量范围 \\
\code{learning_rate / min_lr_ratio} & $10^{-4}/0.1$ & 基础 LR 与最终比例 \\
\code{warmup_steps} & $1000$ & 线性预热后余弦衰减 \\
\code{weight_decay / gradient_clip} & $0.01/1$ & AdamW 衰减与全局 L2 阈值 \\
\code{author_style_loss_weight} & $0.05$ & 作者项权重 \\
\code{precision / preload_to_gpu / workers} & BF16 / False / $8$ & 精度与数据存放方式 \\
\code{val_batches / validate_every} & $16/1000$ & 验证批数与间隔 \\
\code{save_every / log_every} & $500/50$ & 保存与日志间隔 \\
\code{sample_steps / seed} & $32/42$ & 预览 Euler 步数与种子 \\
\bottomrule
\end{tabularx}\end{center}
新训练读取 \code{src/config.py}，AdamW betas 为 $(0.9,0.95)$，全部参数同组衰减。

CUDA 支持 FP32、BF16、FP16 并启用 TF32；CPU 实际为 FP32，只有 FP16 启用 GradScaler。
先反缩放、诊断和裁剪梯度，再更新参数；范数归约非有限时以 float64 重算。
真实异常梯度不更新参数：FP16 降低 scale 并复用同一 batch/时间/噪声，最多尝试 $16$ 次；
BF16/FP32 或非有限损失直接停止，保留成功保存的检查点。

采样时先计算 $C_\star,U$，从 $x_0\sim\mathcal N(0,I)$ 出发，$J$ 步 Euler：
\[
t_j=j/J,\qquad x_{j+1}=x_j+J^{-1}v_\theta(x_j,t_j,C_\star,U),\quad j=0,\ldots,J-1.
\]
只在返回时裁剪 $[-1,1]$，灰度映射由调用方完成。可提供固定 noise 复现实验。
没有额外随机积分量；内容和风格每次采样各编码一次，主干每步仍重算各层 memory LN/K/V，没有 K/V cache。

\section{数据、验证与恢复}

从 \code{DiT-flow/} 运行 \code{python -m src.prepare_data}，必填 GNT 根目录、字体和输出路径。
默认边长 $64$、边距 $4$、字体索引 $0$、最少字符数 $5$、验证比例 $0.1$、种子 $42$。
每个 GNT 文件视为作者，文件 stem 是 ID，要求至少两个文件且 ID 不重复；数值字段小端、标签按 GBK 解码。
保留字体支持的基本 CJK 单字，去除重复、空白与字符数不足作者；损坏记录直接报错。
手写和标准字紧裁白背景、保持宽高比、LANCZOS 缩放居中，无随机图像增强。
只应提供训练集文件，代码不自动识别官方 Test 文件。

按作者划分 train/val，验证人数为 $\min(W-1,\max(1,\operatorname{round}(0.1W)))$（默认比例）。
作者不重叠，字符类别可以重叠；保存 uint8 NPY 和最后写入的版本 $1$ manifest。
训练和验证全部字图预加载一份：CPU 使用 RawArray 供 worker 共享；CUDA 常驻则强制 worker 为 $0$。
episode 独立随机数由种子与全局索引确定，续训从 saved step 乘 batch size 继续索引，验证用 seed 加一的固定 episode。
每字图转为 pixel/127.5-1，参考补白并附掩码和稳定作者索引；不再返回参考标准字。
数据指纹为 manifest 原始字节 SHA-256，不重新核验记录的 NPY/字体内容哈希。
旧预处理数据满足尺寸及最大参考数加一个不同字符时可复用。

验证每隔固定间隔及计划末步进行，使用当前模型双损失与预览，结束恢复 train/eval 状态。
隔离 RNG、使用 seed 加 $10000$ 固定时间/噪声，全部损失算完再采样预览。
\code{--stop-after} 是绝对优化器步数，不改 LR 总日程。训练循环响应 SIGINT/SIGTERM：
完成当前更新后跳过验证并保存 last；强杀或一般异常不保证保存当前步。

检查点版本为 \textbf{3}，仅保存当前模型、AdamW、GradScaler、PyTorch RNG、配置、数据指纹、步数和最优验证值。
版本 $1/2$ 架构不兼容，在写输出前拒绝恢复；使用新目录重新训练。
best 在当前模型验证总目标严格改善时保存；last 按间隔及结束保存，可能落后于内存当前步。
恢复采用保存超参数，运行开关 \code{preload_to_gpu} 读取当前配置，设备、worker 与提前终点可由命令行指定。
新训练拒绝非空结果；续训旧目录配置/指纹须一致，且已有 last 不得比选中恢复点新。
保存先写临时文件、flush/fsync 再替换；I/O 错误最多三次尝试，间隔 $1$、$2$ 秒，不重复训练。
刷新 best 与 last 不是整体事务，存储需额外一份完整检查点空间。

\section{计算规模与验证边界}

参考主要注意力成本为 $O(B_sNL^2d+NZLd)$，内容为 $O(B_cL^2d)$，
$J$ 步主干为 $O(JB_g(2L^2+LNZ)d)$，另需计入投影、FFN、batch 与反向。
实际参考计算按补齐槽数计；PyTorch MHA 后端由环境决定，不保证使用 FlashAttention。
现有测试覆盖双条件结构与梯度、执行顺序、位置编码、掩码/重排、作者配对平均、纯噪声 flow 与缓存采样，
以及数据共享、预处理、诊断、写入失败和 CPU 小模型续训一致性。
模拟设备测试与 CPU 回归不能替代真实 GPU 性能及生成质量实验。
\end{document}
